\documentclass[10pt,a4paper]{amsart}
\usepackage[margin=19mm]{geometry}
\usepackage{amsmath,amssymb,mathtools}
\usepackage[scr=rsfs,cal=euler]{mathalfa}
\usepackage{xfrac}
\usepackage[unicode,hidelinks,bookmarks=false]{hyperref}
\newcommand{\ie}{i.e.\ }
\newcommand{\cf}{cf.\ }
\newcommand{\resp}{respectively\ }
\newcommand{\Subsection}{\subsection}
\providecommand{\nobreakdash}{\nobreak\hbox{-}}
\setlength{\emergencystretch}{2em}
\multlinegap0pt
\usepackage{amssymb}
\usepackage{xcolor}
\usepackage{enumerate}
\usepackage{algorithm}\usepackage{algpseudocode}
\newtheorem{thm}{Theorem}
\newcommand\bfI{\mathbf{I}}
\newcommand\bfa{\mathbf{a}}
\newcommand\bfb{\mathbf{b}}
\newcommand\bff{\mathbf{f}}
\newcommand\bfv{\mathbf{v}}
\newcommand\bfw{\mathbf{w}}
\newcommand\bfx{\mathbf{x}}
\title{The Curse of Black Sigatoka: A Backward Bifurcation Perspective}
\author{Bernard Asamoah Afful, Luis F. Gordillo}
\date{Source: arXiv:2604.26060v1 (CC BY 4.0)}
\begin{document}
\maketitle

\noindent\emph{Source and licence.} The Curse of Black Sigatoka: A Backward Bifurcation Perspective. Version arXiv:2604.26060v1 (CC BY 4.0), distributed by arXiv under the Creative Commons Attribution 4.0 International licence, http://creativecommons.org/licenses/by/4.0/, as recorded in the arXiv licence metadata for this exact version and shown on the abstract page https://arxiv.org/abs/2604.26060v1. Full licence notice: \texttt{licenses/arxiv-2604.26060-CC-BY-4.0.txt}.\par\smallskip

\noindent\textit{Editorial scope.} Counted unit: the article's thm thm (source0.tex lines 409--419). The article's complete original proof of it is delivered with it (source0.tex lines 420--537, one \texttt{proof} environment of 118 lines). No mathematics is written, repaired or re-derived: the statement and the proof are the article's own lines, in the article's own notation, and no retained line is substituted or re-flowed.  Retained uncounted as the source-local context the proof needs: lines 175--188; lines 326--328.  Nothing was omitted from any retained span, so the package carries no omission record.  No retained line contains a citation, so the wrapper carries no bibliography and no bibliography entry.  No retained line contains a figure environment, an image-inclusion command or drawing code, so no image is delivered and the package declares no asset dependency.  Editorial formatting only, in addition: an AMS \texttt{amsart} preamble that restates the article's own declarations and macros used by the retained lines, namely the environments \texttt{thm} on separate counters, together with the standard packages those lines need.  The retained environments print in this document as Theorem 1 in order of appearance, whereas the article prints the same items with the numbering of its own full document.  The complete native source of the article as distributed in the arXiv e-print for this exact version is bundled unchanged as \texttt{sources/199361/source0.tex}.
\begin{align}\label{eqn1}
    \text{\textbf{Host population}}&
    \begin{cases}
        \dfrac{dH}{dt}=\kappa\Lambda - \beta(h,T)(\theta S+\psi\frac{A}{N})H - \mu H,\\[10pt]
	\dfrac{dR}{dt}=(1-\kappa)\Lambda - (1-\delta)\beta(h,T)(\theta S+\psi\frac{A}{N})R - \mu R,\\[10pt]
	\dfrac{dE}{dt}=\beta(h,T)(\theta S+\psi \frac{A}{N})(H+(1-\delta)R) - (\gamma+\mu)E,\\[10pt]
	\dfrac{dI}{dt}=\gamma E - (\mu+\rho)I.
    \end{cases}\\[10pt]
    \text{\textbf{Pathogen population}}&
    \begin{cases}
        \dfrac{dA}{dt}=\eta I - (\mu_P+\rho)A,\\[10pt]
	\dfrac{dS}{dt}=\alpha g(I)I - (\mu_P+\rho)S.
    \end{cases}
\end{align}

\begin{equation}\label{eqn5}
    \mathcal{R}_0=\frac{\beta(h,T)\psi\eta\gamma[\delta\kappa+(1-\delta)]}{(\gamma+\mu)(\mu_P+\rho)(\mu+\rho)}.
\end{equation}

\begin{thm}
    Let us denote $\bfx=(H,R,E,I,A,S)$ and $\bff=d\bfx/dt$. The bifurcation coefficients $\bfa$ and $\bfb$ are defined as
    \begin{equation*}
        \bfa=\sum_{k,i,j=1}^{6}v_kw_iw_j\frac{\partial^2f_k(\mathcal{E}^0,\beta_0^*)}{\partial x_i\partial x_j},\qquad
        \bfb=\sum_{k,i=1}^{6}v_kw_i\frac{\partial^2f_k(\mathcal{E}^0,\beta_0^*)}{\partial x_i\partial\beta_0^*}.
    \end{equation*}
    Then, the BSD model \eqref{eqn1} undergoes a backward bifurcation at $\beta_0=\beta_0^*$, i.e. $\mathcal{R}_0=1$, when $\bfa>0$, $\bfb>0$ and $\lambda>\lambda^{\text{crit}}$, where
    \begin{equation}
        \lambda^{\text{crit}}=\frac{\beta_0^*\Pi\mu\psi^2\eta^2[(1-\delta)+\delta\kappa(2+\delta)]}{\alpha\theta\Lambda^2(\mu_P+\rho)[\delta\kappa+(1-\delta)]}.
    \end{equation}
\end{thm}

\begin{proof}
    Setting $\mathcal{R}_0=1$ and solving for $\beta_0^*$ in \eqref{eqn5} gives;
    \begin{equation}\label{eqn7}
        \beta_0^*=\frac{(\gamma+\mu)(\mu_P+\rho)(\mu+\rho)}{\Pi\psi\eta\gamma[\delta\kappa+(1-\delta)]},
    \end{equation}
    where $\Pi=\frac{h}{h+K_h}\exp{\left(-\frac{(T-\hat{T})^2}{2\sigma_T^2}\right)}$. The Jacobian matrix of the BSD model \eqref{eqn1} evaluated at the BSD-free equilibrium, $\mathcal{E}^0$ with $\beta_0=\beta_0^*$ is
    \begin{equation}\label{eqn8}
        \mathcal{J}(\mathcal{E}^0,\beta_0^*)=
        \begin{bmatrix}
            -\mu&0&0&0&-\beta_0^*\Pi\psi\kappa&-\frac{\beta_0^*\Pi\theta\kappa\Lambda}{\mu}\\
            0&-\mu&0&0&-\beta_0^*\Pi\psi(1-\delta)(1-\kappa)&-\frac{\beta_0^*\Pi\theta\Lambda(1-\delta)(1-\kappa)}{\mu}\\
            0&0&-(\gamma+\mu)&0&\beta_0^*\Pi\psi[\delta\kappa+(1-\delta)]&\frac{\beta_0^*\Pi\theta\Lambda[\delta\kappa+(1-\delta)]}{\mu}\\
            0&0&\gamma&-(\mu+\rho)&0&0\\
            0&0&0&\eta&-(\mu_P+\rho)&0\\
            0&0&0&0&0&-(\mu_P+\rho)
        \end{bmatrix}.
    \end{equation}
    This matrix has $\lambda_1=\lambda_2=-\mu$, and $\lambda_6=-(\mu_P+\rho)$ as eigenvalues. The remaining eigenvalues are obtained from the sub-matrix
    \begin{equation}\label{eqn9}
        \hat{\mathcal{J}}(\mathcal{E}^0,\beta_0^*)=
        \begin{bmatrix}
            -(\gamma+\mu)&0&\beta_0^*\Pi\psi[\delta\kappa+(1-\delta)]\\
            \gamma&-(\mu+\rho)&0\\
            0&\eta&-(\mu_P+\rho)
        \end{bmatrix},
    \end{equation}
    for which the associated characteristic polynomial at $\chi$ is
    \begin{equation}
        |\chi\bfI-\hat{\mathcal{J}}(\mathcal{E}^0,\beta_0^*)|=(\chi+\gamma+\mu)(\chi+\mu+\rho)(\chi+\mu_P+\rho)-\beta_0^*\Pi\psi\gamma\eta[\delta\kappa+(1-\delta)].
    \end{equation}
    At $\chi=0$, there is an eigenvalue when
    \begin{align*}
        (\gamma+\mu)(\mu+\rho)(\mu_P+\rho)-\beta_0^*\Pi\psi\gamma\eta[\delta\kappa+(1-\delta)]&=0,\\
        \frac{\beta_0^*\Pi\psi\gamma\eta[\delta\kappa+(1-\delta)]}{(\gamma+\mu)(\mu+\rho)(\mu_P+\rho)}&=1,\\
        \mathcal{R}_0&=1.
    \end{align*}
    Therefore, at $\beta_0=\beta_0^*$ that is $\mathcal{R}_0=1$, the Jacobian matrix associated with the BSD model \eqref{eqn1}  evaluated at $\mathcal{E}^0$ has a simple zero eigenvalue (with all other eigenvalues having negative real part). Thus, the Jacobian matrix has a right eigenvector given by $\bfw=(w_1,w_2,w_3,w_4,w_5,w_6)$ and a left eigenvector given by $\bfv=(v_1,v_2,v_3,v_4,v_5,v_6)$ corresponding to the zero eigenvalue. A right eigenvector satisfies
    \begin{equation}\label{eqn10}
        \begin{bmatrix}
            -\mu&0&0&0&-\beta_0^*\Pi\psi\kappa&-\frac{\beta_0^*\Pi\theta\kappa\Lambda}{\mu}\\
            0&-\mu&0&0&-\beta_0^*\Pi\psi(1-\delta)(1-\kappa)&-\frac{\beta_0^*\Pi\theta\Lambda(1-\delta)(1-\kappa)}{\mu}\\
            0&0&-(\gamma+\mu)&0&\beta_0^*\Pi\psi[\delta\kappa+(1-\delta)]&\frac{\beta_0^*\Pi\theta\Lambda[\delta\kappa+(1-\delta)]}{\mu}\\
            0&0&\gamma&-(\mu+\rho)&0&0\\
            0&0&0&\eta&-(\mu_P+\rho)&0\\
            0&0&0&0&0&-(\mu_P+\rho)
        \end{bmatrix}
        \begin{bmatrix}
            w_1\\w_2\\w_3\\w_4\\w_5\\w_6
        \end{bmatrix}
        =
        \begin{bmatrix}
            0\\0\\0\\0\\0\\0
        \end{bmatrix}
    \end{equation}
    Solving for $w_1$, $w_2$, $w_3$, and $w_5$ in terms of $w_4$ yields
    \begin{align*}
        w_1&=-\frac{\beta_0^*\Pi\psi\kappa\eta}{\mu(\mu_P+\rho)}w_4,\;\;
        w_2=-\frac{\beta_0^*\Pi\psi\eta(1-\delta)(1-\kappa)}{\mu(\mu_P+\rho)}w_4,\;\;
        w_3=\frac{\mu+\rho}{\gamma}w_4,\\
        w_4&>0,\;\;
        w_5=\frac{\eta}{\mu_P+\rho}w_4,\;\;
        w_6=0.
    \end{align*}
    A left eigenvector satisfies
    \begin{equation}\label{eqn11}
        \begin{bmatrix}
            -\mu&0&0&0&0&0\\
            0&-\mu&0&0&0&0\\
            0&0&-(\gamma+\mu)&\gamma&0&0\\
            0&0&0&-(\mu+\rho)&\eta&0\\
            -\beta_0^*\Pi\psi\kappa&-\beta_0^*\Pi\psi(1-\delta)(1-\kappa)&\beta_0^*\Pi\psi[\delta\kappa+(1-\delta)]&0&-(\mu_P+\rho)&0\\
            -\frac{\beta_0^*\Pi\theta\kappa\Lambda}{\mu}&-\frac{\beta_0^*\Pi\theta\Lambda(1-\delta)(1-\kappa)}{\mu}&\frac{\beta_0^*\Pi\theta\Lambda[\delta\kappa+(1-\delta)]}{\mu}&0&0&-(\mu_P+\rho)
        \end{bmatrix}
        \begin{bmatrix}
            v_1\\v_2\\v_3\\v_4\\v_5\\v_6
        \end{bmatrix}
        =
        \begin{bmatrix}
            0\\0\\0\\0\\0\\0
        \end{bmatrix}
    \end{equation}
    Solving for $v_4$, $v_5$, and $v_6$ in terms of $v_3$ gives
    \begin{align*}
        v_1&=0,\;
        v_2=0,\;
        v_3>0,\;\;
        v_4=\frac{\gamma+\mu}{\gamma}v_3,\;\;
        v_5=\frac{(\mu+\rho)(\gamma+\mu)}{\eta\gamma}v_3,\;\;
        v_6=\frac{\beta_0^*\Pi\theta\Lambda[\delta\kappa+(1-\delta)]}{\mu(\mu_P+\rho)}v_3.
    \end{align*}
    Since $v_1$ and $v_2$ are zero, we only compute the second-order partial derivatives of $f_3$, $f_4$, $f_5$,  and $f_6$,
    \begin{align*}
        \frac{\partial^2f_3(\mathcal{E}^0,\beta_0^*)}{\partial H\partial S}&=\frac{\partial^2f_3(\mathcal{E}^0,\beta_0^*)}{\partial S\partial H}=\beta_0^*\Pi\theta,\\
        \frac{\partial^2f_3(\mathcal{E}^0,\beta_0^*)}{\partial H\partial A}&=\frac{\partial^2f_3(\mathcal{E}^0,\beta_0^*)}{\partial A\partial H}=\frac{\beta_0^*\Pi\psi\mu}{\Lambda},\\
        \frac{\partial^2f_3(\mathcal{E}^0,\beta_0^*)}{\partial R\partial S}&=\frac{\partial^2f_3(\mathcal{E}^0,\beta_0^*)}{\partial S\partial R}=\beta_0^*\Pi\theta(1-\delta),\\
        \frac{\partial^2f_3(\mathcal{E}^0,\beta_0^*)}{\partial R\partial A}&=\frac{\partial^2f_3(\mathcal{E}^0,\beta_0^*)}{\partial A\partial R}=\frac{\beta_0^*\Pi\psi\mu(1-\delta)}{\Lambda},\\
        \frac{\partial^2f_6(\mathcal{E}^0,\beta_0^*)}{\partial I^2}&=2\alpha\lambda,\\
        \frac{\partial^2f_3(\mathcal{E}^0,\beta_0^*)}{\partial A\partial\beta_0^*}&=\Pi\psi\mu[\delta\kappa+(1-\delta)],\\
        \frac{\partial^2f_3(\mathcal{E}^0,\beta_0^*)}{\partial S\partial\beta_0^*}&=\frac{\Pi\theta\Lambda[\delta\kappa+(1-\delta)]}{\mu}.
    \end{align*}
    Using these expressions in the definitions of $\bfa$ gives 
    \begin{equation}
        \bfa=2v_3w_5\frac{\beta_0^*\Pi\psi\mu}{\Lambda}\left[w_1+(1-\delta)w_2\right]+2v_6w_4^2\alpha\lambda.
    \end{equation}
    Substituting the values of $w_1$, $w_2$, $w_5$, and setting $w_4=1$ yields
    \begin{equation}
        \bfa=\frac{2v_3\alpha\lambda\beta_0^*\Pi\theta\Lambda[\delta\kappa+(1-\delta)]}{\mu(\mu_P+\rho)}-\frac{2v_3(\beta_0^*\Pi\psi\eta)^2[(1-\delta)+\delta\kappa(2+\delta)]}{\Lambda(\mu_P+\rho)}.
    \end{equation}
    Therefore, if 
    \begin{equation}\label{eqn12}
        \lambda>\underbrace{\frac{\beta_0^*\Pi\mu\psi^2\eta^2[(1-\delta)+\delta\kappa(2+\delta)]}{\alpha\theta\Lambda^2(\mu_P+\rho)[\delta\kappa+(1-\delta)]}}_{\lambda^{\text{crit}}}
    \end{equation}
    then $\bfa>0$. Moreover, 
    \begin{equation}
        \bfb=v_3w_5\Pi\psi[\delta\kappa+(1-\delta)]=\frac{\Pi\psi\eta[\delta\kappa+(1-\delta)]}{\mu_P+\rho}v_3>0.
    \end{equation}
    By applying Theorem 4.1 of Castillo-Chavez and Song, we conclude that the BSD model \eqref{eqn1} undergoes a backward bifurcation. 
    \end{proof}

\end{document}
